% !TEX program = xelatex
% 이전 gate 실험(실험11 MoE M0)과 이번 실험(실험16 Gated Late Fusion)의 차이 — 보고용 한 장 요약.
% 수치 출처: outputs/moe_gate_m0_radbert/results.json (M0, RadBERT), archive/before-cleanup-2026-09-30 태그의
%            experiments/실험11_MoE/RESULTS_ALL.md (M0-tuned, M1, 게이트 표준편차),
%            outputs/gmu_nested/gated_nested/results.json · report_v2/summary.json (실험16).
\documentclass[tikz,border=8pt]{standalone}
\usepackage{fontspec}
\setmainfont{Noto Sans CJK KR}
\setsansfont{Noto Sans CJK KR}
\usepackage{amsmath,amssymb}
\usetikzlibrary{positioning,arrows.meta,calc}

\definecolor{navy}{HTML}{2F3B52}
\definecolor{ink}{HTML}{1A1A1A}
\definecolor{ink2}{HTML}{5A5E66}
\definecolor{boxblue}{HTML}{CFE0F7}
\definecolor{boxpink}{HTML}{F6CDD0}
\definecolor{boxyellow}{HTML}{FDEEB8}
\definecolor{boxgray}{HTML}{E9E9E9}
\definecolor{bad}{HTML}{C9545C}
\definecolor{good}{HTML}{3E6FB0}

\tikzset{
  every node/.style={text=ink},
  cell/.style={draw=navy, line width=0.6pt, rounded corners=3pt, fill=#1, align=left,
               text width=9.4cm, minimum height=2.05cm, inner xsep=7pt, inner ysep=5pt, font=\small},
  rowlab/.style={draw=navy, line width=0.6pt, rounded corners=3pt, fill=boxgray, align=center,
                 text width=2.5cm, minimum height=2.05cm, font=\small\bfseries},
  head/.style={font=\normalsize\bfseries, align=center, text width=9.4cm},
  arr/.style={-{Latex[length=2.6mm,width=2mm]}, draw=navy, line width=1pt},
  note/.style={font=\scriptsize, text=ink2, align=left},
}

\begin{document}
\begin{tikzpicture}[x=1cm,y=1cm]

% ---------------- 제목 ----------------
\node[anchor=west, font=\Large\bfseries] at (0,0.9) {이전 gate 실험이 낮았던 이유와 이번 실험에서 바꾼 점};
\node[anchor=west, font=\small, text=ink2] at (0,0.25)
  {비교 대상: 위험점수를 섞는 gate — 이전 실험11 (MoE 게이트 M0) → 이번 실험16 (Gated Late Fusion)};

% ---------------- 열 제목 ----------------
\node[head, text=bad]  at (8.2,-0.55)  {이전 실험 (실험11, M0)};
\node[head, text=good] at (19.25,-0.55) {이번 실험 (실험16, Gated Late Fusion)};

% ---------------- 3개 행 ----------------
\def\rA{-1.85} \def\rB{-4.1} \def\rC{-6.35}

\node[rowlab] (l1) at (1.3,\rA) {① 게이트에\\ 넣은 입력};
\node[cell=boxpink] (b1) at (8.2,\rA)
  {모달리티별 임베딩을 PCA로 줄여 \textbf{24차원}을 입력\\
   → 학습 환자 171명에 비해 입력이 많아 과적합\\
   {\footnotesize\color{ink2}(attention 게이트 M1은 학습 값 약 1.1만 개, OS 차이 −0.12로 중단)}};
\node[cell=boxblue] (a1) at (19.25,\rA)
  {두 모델의 위험점수 \textbf{2개}($z_{\text{tab}}, z_{\text{img}}$)만 입력\\
   학습하는 값은 \textbf{6개}($\beta$ 2개, $c$ 4개)\\
   {\footnotesize\color{ink2}게이트 식: $g_m = 2\,\sigma(c_m^{\top}z)$}};

\node[rowlab] (l2) at (1.3,\rB) {② 게이트를\\ 끄는 장치};
\node[cell=boxpink] (b2) at (8.2,\rB)
  {\textbf{없음}\\
   도움이 안 되는데도 환자마다 게이트 값이 크게 달라짐\\
   {\footnotesize\color{ink2}(영상 게이트 표준편차 약 0.3, 시험 환자에서는 성능 하락)}};
\node[cell=boxblue] (a2) at (19.25,\rB)
  {$c$에 벌점 $\lambda\lVert c\rVert^2$을 주고, $\lambda$는 inner CV로 선택\\
   $c=0$이면 기존 late fusion과 \textbf{같은 식}이 됨\\
   {\footnotesize\color{ink2}실제로 OS 5개 fold 중 2개, PFS 5개 중 2개에서 게이트가 꺼짐}};

\node[rowlab] (l3) at (1.3,\rC) {③ 게이트 학습에\\ 쓴 점수};
\node[cell=boxpink] (b3) at (8.2,\rC)
  {각 모델이 \textbf{학습에 쓴 환자를 다시 채점한 점수}로 학습\\
   → 기준 모델 값부터 높게 나옴 (OS 0.730)\\
   {\footnotesize\color{ink2}그래서 기존 결과(0.722)와 값을 직접 비교할 수 없었음}};
\node[cell=boxblue] (a3) at (19.25,\rC)
  {inner CV에서 나온 \textbf{OOF 점수}(학습에 안 쓴 환자 점수)로 학습\\
   시험 환자는 학습·선택에 전혀 쓰지 않음 (nested CV)};

\foreach \i in {1,2,3} { \draw[arr] (b\i.east) -- (a\i.west); }

% ---------------- 아래 왼쪽: Δ 막대그래프 ----------------
\begin{scope}[shift={(0.6,-10.4)}]
  \node[anchor=west, font=\normalsize\bfseries] at (-0.4,2.55) {게이트를 넣었을 때 C-index 차이 (Δ)};
  \node[anchor=west, note] at (-0.4,2.05) {Δ = 게이트 모델 − 같은 실험의 기준 모델 (5-fold 평균)};
  % 축: 0.01 → 1.0 cm
  \def\s{100}
  \draw[navy!40, line width=0.4pt] (0,{0.01*\s}) -- (11.5,{0.01*\s});
  \draw[navy!40, line width=0.4pt] (0,{-0.01*\s}) -- (11.5,{-0.01*\s});
  \draw[navy!40, line width=0.4pt] (0,{-0.02*\s}) -- (11.5,{-0.02*\s});
  \draw[navy, line width=0.8pt] (0,0) -- (11.5,0);
  \draw[navy, line width=0.6pt] (0,{-0.027*\s}) -- (0,{0.013*\s});
  \foreach \v/\t in {0.01/+0.01, 0/0, -0.01/$-$0.01, -0.02/$-$0.02}
    \node[anchor=east, font=\scriptsize, text=ink2] at (-0.08,{\v*\s}) {\t};
  % 막대: OS
  \fill[bad!75]  (1.2,0) rectangle (2.8,{-0.0238*\s});
  \fill[good!80] (3.2,0) rectangle (4.8,{0.0014*\s});
  \node[font=\footnotesize\bfseries, text=bad]  at (2.0,{-0.0238*\s-0.3}) {$-$0.024};
  \node[font=\footnotesize\bfseries, text=good] at (4.0,{0.0014*\s+0.3}) {+0.001};
  \node[font=\small\bfseries] at (3.0,1.45) {OS};
  % 막대: PFS
  \fill[bad!75]  (6.9,0) rectangle (8.5,{-0.0073*\s});
  \fill[good!80] (8.9,0) rectangle (10.5,{-0.0002*\s});
  \node[font=\footnotesize\bfseries, text=bad]  at (7.7,{-0.0073*\s-0.3}) {$-$0.007};
  \node[font=\footnotesize\bfseries, text=good] at (9.7,0.3) {0.000};
  \node[font=\small\bfseries] at (8.7,1.45) {PFS};
  % 범례
  \fill[bad!75]  (1.2,-3.35) rectangle +(0.35,0.3); \node[anchor=west, font=\scriptsize] at (1.6,-3.2) {이전 M0 (RadBERT) — 기준: 같은 실험의 상수 게이트};
  \fill[good!80] (1.2,-3.85) rectangle +(0.35,0.3); \node[anchor=west, font=\scriptsize] at (1.6,-3.7) {이번 Gated Late Fusion — 기준: nested late fusion};
  \node[anchor=west, note] at (-0.4,-4.3) {두 실험은 평가 방식이 달라 C-index 값 자체가 아니라 차이(Δ)만 비교함};
\end{scope}

% ---------------- 아래 오른쪽: 결과 정리 ----------------
\node[draw=navy, line width=0.6pt, rounded corners=3pt, fill=boxyellow!55, anchor=north east,
      text width=11.2cm, align=left, inner xsep=9pt, inner ysep=8pt, font=\small] at (24.2,-7.6)
  {\textbf{이번 실험 C-index (5-fold 평균)}\\[2pt]
   \begin{tabular}{@{}lcc@{}}
     & OS & PFS\\ \hline
     Late fusion (기준) & 0.723 & 0.674\\
     Gated late fusion  & 0.725 & 0.674\\
   \end{tabular}\\[8pt]
   \textbf{정리}\\[2pt]
   • 이전에는 게이트를 넣으면 성능이 떨어졌지만(OS $-$0.024), 이번에는 떨어지지 않았다.\\[2pt]
   • 다만 late fusion보다 높아지지도 않았다(OS +0.001, PFS 0.000).\\[2pt]
   • 세 가지를 한 번에 바꿔서 어느 쪽 영향이 가장 컸는지는 따로 확인하지 않았다.
     이전 실험에서도 입력을 줄이고 규제를 건 버전(M0-tuned)은 차이가 OS $-$0.003까지 줄었다.\\[2pt]
   • 따라서 이번 데이터(238명)에서는 환자별 게이트로 얻는 추가 이득은 확인되지 않았다.};

\end{tikzpicture}
\end{document}
